\subsection{The output of the application, panel by panel}
The water of Chapter 1 is not only a set of equations: in SWAFL it is a set of pictures that change with the settings of the user. The layer \emph{Soil hydraulic drivers} shows how the retention curve and the conductivity of the soil move when the tillage, the consolidation and the texture change, and the figures below are its panels for the coarse loam.

The first panel ranks the ten soils by the change of a hydraulic property (the saturated conductivity or the water held) when one driver is moved from its present value. The bars are the sensitivity of the property to the driver; the scatter on the right places the soils on the plane of the two drivers. It is a direct reading of the van Genuchten parameters of this chapter as functions of the pore structure.

\appfig{v7_0}{0.96}{What drives change in soil hydraulic properties: the response of the soil to tillage, consolidation rate and organic matter, soil by soil.}

The graphs follow the drainage against the consolidation rate and the water content against the suction for conventional, reduced and no-till systems; clicking a point applies that setting to the whole program, so the effect of a change on the retention curve can be seen at once.

\appfig{v7_1}{0.96}{Graphs of the drivers at your settings: drainage and water content against the consolidation rate, for the three tillage systems.}

The pictures are sections through the topsoil of the ensemble runs after 0, 50, 100 and 200 days of settling. The dark cracks and macropores close as the soil consolidates, and the pore index printed under each section (0.92 and 0.67 for two states) is the number that enters the conductivity of the layer.

\appfig{v7_2}{0.96}{Real soil profiles: how tillage and settling change the soil, layer by layer, in the pictures of the model.}

